\documentclass[11pt]{article}

% Packages
\usepackage[margin=1in]{geometry}
\usepackage{setspace}
\usepackage{parskip}
\usepackage{times}
\usepackage{fancyhdr}

% Header setup
\pagestyle{fancy}
\fancyhf{}
\chead{\small Scott Nguyen --- Statement of Purpose --- Stanford Aeronautics and Astronautics Ph.D.}
\renewcommand{\headrulewidth}{0pt}

% No indent
\setlength{\parindent}{0pt}
\onehalfspacing

\begin{document}
	
	My interest in autonomous spacecraft Guidance, Navigation, and Control (GNC) began during my undergraduate aerospace engineering coursework, when I first encountered the challenge of how a spacecraft determines and controls its state with only limited sensing. The structure of nonlinear dynamics, filtering, and estimation aligned naturally with how I reason about problems, and I pursued advanced courses, research, and internships to deepen these skills. Over time, I realized that my strongest abilities—analytical modeling, coding, and systems-level thinking—map directly onto the growing need for resilient autonomy in space systems.
	
	My academic preparation reflects a deliberate effort to build a strong foundation for research in estimation and autonomous navigation. I completed an M.S. in Aerospace Engineering at the University of Illinois Urbana–Champaign, where I focused on orbital mechanics, estimation methods, and spacecraft dynamics. I am now completing a second M.S. in Electrical/Space Systems Engineering at the University of New Mexico to strengthen my background in nonlinear filtering, stochastic modeling, and sensor fusion. Across these programs, I gained experience with the Extended and Unscented Kalman Filters, covariance propagation, Monte Carlo methods, and high-fidelity modeling that incorporates perturbations, atmospheric effects, and relative orbital motion. This academic path has prepared me for research that integrates multi-sensor estimation, optical navigation, and autonomous GNC for multi-satellite and deep-space missions.
	
	My professional experiences reinforced this trajectory. At the Space Dynamics Laboratory, I worked on angles-only orbit estimation and learned firsthand the challenges of observability, sparse sensing, and estimator tuning under real measurement constraints. At Blue Canyon Technologies, I supported hardware-in-the-loop testing of star trackers, inertial measurement units, reaction wheels, and solar array drive assemblies, and assisted in flight-software verification using COSMOS and C-based S-functions. This gave me exposure to how GNC algorithms behave on real hardware and how sensor characteristics shape achievable navigation performance.
	
	At Varda Space Industries, I implemented an EKF for re-entry capsule state estimation, built Monte Carlo simulation tools for descent and recovery analysis, and evaluated gravity models to improve robustness. At Blue Origin, I designed hybrid Active Disturbance Rejection Control–Sliding Mode Control architectures and developed closed-loop simulation frameworks to evaluate attitude-control robustness. Most recently, at MIT Lincoln Laboratory, I have contributed to research on radio-frequency–based state estimation, probabilistic detection algorithms, and S-band communication and ranging architectures for resilient navigation in contested environments. Together, these roles broadened my perspective on spacecraft autonomy—from low-Earth orbit to re-entry to distributed satellite systems—and helped refine my research interests toward estimation and decision-making in complex, uncertain regimes.
	
	I am applying to the Stanford Aeronautics and Astronautics Ph.D. program because it offers the ideal intellectual environment to pursue advanced research in autonomous navigation, multi-satellite systems, and vision-based GNC. Stanford’s strengths in nonlinear dynamics, multi-agent systems, and distributed space architectures align precisely with the problems I hope to investigate. In particular, I am strongly motivated by the research conducted in the Space Rendezvous Laboratory (SLAB). The lab’s work in angles-only navigation, formation flying, optical sensing, machine-learning–enabled pose estimation, and hardware-in-the-loop validation represents the exact intersection of topics that motivate my academic and professional goals. SLAB's approach—integrating fundamental theory, algorithm development, and mission-oriented demonstration—is the research philosophy I want to be part of.
	
	My long-term research interests center on developing navigation and estimation architectures that allow spacecraft to operate with minimal dependence on the Deep Space Network or Global Positioning System. I am particularly interested in: (1) optical and RF-based relative navigation for swarms and formation-flying spacecraft; (2) image-based pose estimation and terrain-relative navigation using machine learning and uncertainty-aware filtering; (3) nonlinear guidance and multi-sensor fusion methods for missions operating in cislunar and multi-body regimes; and (4) trajectory design and propagation in the Circular Restricted Three-Body Problem with stochastic uncertainties. These areas map directly onto the lab’s thrusts in multi-agent GNC, distributed space systems, and vision-based navigation, and I believe I can contribute meaningfully to ongoing projects while proposing new extensions related to autonomy in weak-signal, degraded-sensing environments.
	
	My future career plans are to work as a research engineer advancing autonomous GNC for space missions. I hope to contribute to technologies that enable multi-satellite architectures, in-space servicing, deep-space exploration, and distributed sensing platforms. A Ph.D. from Stanford would provide the rigorous theoretical training, interdisciplinary collaboration, and mission-focused research environment needed to build such a career. I am particularly drawn to the program’s integration of fundamental dynamics with real mission testbeds, its strong ties to NASA and industry, and its culture of interdisciplinary research across sensing, control, machine learning, and astrodynamics.
	
	For these reasons, I am eager to pursue doctoral study at Stanford and contribute to the Aeronautics and Astronautics community. The program’s depth in dynamics, navigation, and multi-satellite systems makes it the ideal place to develop the skills I need to help build the next generation of autonomous space missions.
	
\end{document}
